\section{Použité technologie}
\subsection{Unreal Engine}
Prostředí i agent se vytváří a používají v Unreal Engine 5.2, který zajišťuje mapu i herní mechanismy s moderními renderovacími funkcemi v reálném čase. Některé z použitých prvků ve hře (hlavně modely a animace) pocházejí z FABu — nové online platformy od Epic Games, která spojuje starý Unreal Marketplace, Sketchfab, Quixel a ArtStation v jednotné řešení.

Unreal Engine (zkráceně UE)\cite{Unreal2025}, vyvíjený společností Epic Games, vzniká z více než dvacetileté evoluce začínající v roce 1998 a dnes představuje špičkový engine pro tvorbu her, filmů i vizualizací. Jeho pružnost, nastavitelnost a výkon si získávají uznání vývojářů po celém světě.

Mezi výrazné přednosti UE patří systém Blueprint vizuálního skriptování — nadstavba nad C++, která umožňuje vytvářet herní logiku, interakce a AI pomocí uzlů a grafů místo klasického kódu v C++. Díky tomu je vývoj her přístupnější i kreativcům bez programátorského zázemí. Když jsou uzly přehledně poskládány, může být práce s nimi dokonce rychlejší pro prototypování a nasazení systémů než použití čistého C++.

Značnou oblibu nachází spíše UE5 než UE obecně — a to nejen mezi herními vývojáři, ale i ve filmové produkci, architektonické vizualizaci, produktovém designu nebo v průmyslu a tréninkových simulacích. Filmové týmy využívají UE5 pro virtuální natáčení v simulovaných prostředích, architekti jej používají v interaktivních prezentacích klientům a v medicíně či obraně slouží k tréninku zaměstnanců ve vysoce realistických scénářích.

\newpage

Důvody volby Unreal Engine oproti alternativám jako Unity:
\begin{enumerate}
	\item C++:
	\begin{itemize}
		\item C++ je jedním ze dvou hlavních jazyků v Unreal Enginu a umožňuje velmi rychlý kód, což je klíčové pro aplikace běžící v reálném čase, jako jsou hry. Navíc je v něm největší zkušenost vývojového týmu (v porovnání s C\#, který je standardní v Unity).
	\end{itemize} 
	\item Programovací nadstavba Blueprint:
	\begin{itemize}
		\item Blueprint je vizuální nadstavba nad C++, která výrazně zjednodušuje tvorbu herní logiky, interakcí a AI. Díky ní mohou tvořit i uživatelé s omezenou znalostí programování — někteří pracovníci s minimálními schopnostmi dosahují produktivity vyšší než při čistém C++.
	\end{itemize} 
	\item FAB:
	\begin{itemize}
		\item Platforma FAB, kterou provozuje Epic Games, sjednocuje dřívější tržiště Unreal Marketplace, Sketchfab, Quixel a ArtStation. Díky tomu poskytuje široké spektrum assetů (modely, animace…), což výrazně urychluje vývoj bez nutnosti tvořit každou položku od nuly.
	\end{itemize}
	\item Pluginy:
	\begin{itemize}
		\item Unreal Engine umožňuje snadné rozšíření funkcionality prostřednictvím pluginů. Ty lze instalovat přímo z oficiálního obchodu, nebo je přidávat ručně do instalační složky UE — flexibilní řešení pro specifické funkce.
	\end{itemize}
	\item 3D architektura:
	\begin{itemize}
		\item UE je navržen od základu pro komplexní 3D tvorbu. V Unity nebo Godotu lze sice vyvíjet 3D aplikace, ale těchto nástrojů se častěji využívá pro 2D či 2.5D projekty kvůli nižší náročnosti.
	\end{itemize}
\end{enumerate}

Nevýhody Unreal Enginu oproti Unity či Godotu:
\begin{enumerate}
	\item Menší komunita:
	\begin{itemize}
		\item Unity Unity má širší uživatelskou základnu, což zajišťuje větší dostupnost řešení specifických problémů. Godot je v tomto ohledu srovnatelný s Unreal Enginem.
	\end{itemize} 
	\item Vyšší hardwarové nároky:
	\begin{itemize}
		\item Unreal Engine 5 pro optimální výkon obvykle vyžaduje 32–64GB RAM, moderní vícejádrový procesor a grafickou kartu minimálně RTX4080, ideálně RTX4090. Epic Games standardně využívá pracovní stanice s 128–256GB RAM, procesory AMD Threadripper a grafikami RTX4080 pro efektivní zpracování rozsáhlých scén. Naproti tomu Unity běží pohodlně i na konfiguraci se 6jádrovým CPU, 16GB RAM a grafickou kartou GTX1060.
	\end{itemize} 
\end{enumerate}

\begin{figure}[h!t]

\includegraphics[scale=0.5]{2-1.png}
\centering
\caption{Ikony použitých technologií}
\label{obrIkon}
\end{figure}

\subsection{Visual Studio}
Visual Studio je používáno jako hlavní vývojové prostředí pro programování v C++ v rámci tohoto projektu. Unreal Engine 5 vyžaduje Visual Studio pro jakoukoli práci s C++, protože jde o úzce integrovaný nástroj — bez něj nelze kompilovat ani spouštět čistě C++ moduly.

Visual Studio je komplexní integrované vývojové prostředí (IDE) vytvořené společností Microsoft. Ačkoli podporuje více programovacích jazyků, při práci s Unreal Enginem se využívá zejména pro C++. Nabízí třeba správu projektů, inteligentní doplňování kódu (IntelliSense), ladění a analýzu paměti.

Při používání Unreal Enginu Visual Studio automaticky stahuje a konfiguruje projektové struktury a závislosti. Využívá build systémy jako MSBuild nebo CMake a úzce spolupracuje se soubory Unreal projektu — díky tomu se nové C++ třídy a funkce správně sestavují a propojují s editorem.

Visual Studio umožňuje ladění v reálném čase, profilování výkonu a sledování využití paměti — klíčové nástroje při vývoji rozsáhlých projektů v UE5.

Při vývoji na Windows je Visual Studio pro Unreal prakticky nezbytné. Umožňuje efektivní práci s C++, řízení kompilace a hladký vývojářský proces při programování na úrovni enginu.

\newpage

\subsection{NevarokML}
NevarokML\cite{Nevarok2025} představuje pokročilý plugin pro Unreal Engine 5, který umožňuje přímou integraci strojového učení — konkrétně posilovaného učení (reinforcement learning, RL) — do herních nebo simulačních projektů. Plugin je založen na knihovně Stable‑Baselines3 a zpřístupňuje oblíbené RL algoritmy jako PPO, A2C, DDPG, DQN, SAC a TD3, které agentům umožňují učit se optimálnímu chování metodou pokus–omyl.

NevarokML podporuje jak jednoagentní, tak víceagentní interakce a nabízí možnost paralelního tréninku napříč odlišnými prostředími. To významně zrychluje učení a pomáhá agentům lépe generalizovat v různých scénářích.

Architektura NevarokML funguje na principu klient–server: Unreal Engine zde působí jako serverová část, která simuluje prostředí a chování agentů, zatímco pythonový klient provádí strojové učení pomocí algoritmů ze Stable‑Baselines3.

Komunikace mezi Unreal Enginem a Python klientem probíhá přes síť nebo sockety s využitím JSON protokolů, které přenášejí informace o stavu prostředí (pozorování, odměny, akce). Jakmile se model natrénuje, výsledné váhy se převedou do formátu ONNX a nasadí se jako objekt typu NNEModelData přímo do Unreal Enginu skrze jeho systém Neural Network Engine.

ONNX\cite{ONNX2025} (Open Neural Network Exchange) představuje otevřený formát a ekosystém navržený tak, aby umožnil snadný přenos již natrénovaných modelů hlubokého učení mezi různými vývojovými frameworky, platformami a typy hardwaru — bez potřeby složité konverze nebo zdlouhavých procesů.

Neural Network Engine (NNE)\cite{NNE2025} je plugin určený pro provádění inferencí neuronových sítí v reálném čase přímo v prostředí Unreal Engine. Plugin využívá formát ONNX, který umožňuje spouštění modelů exportovaných z běžně používaných frameworků strojového učení, jako jsou PyTorch, TensorFlow nebo MXNet, ve formátu .onnx bez nutnosti dalšího převodu.

\begin{enumerate}
\item Vytvoření prostředí:
\begin{itemize}
\item Práce začíná vytvořením Blueprint třídy děděné z NevarokMLEnviroment, kde se definují základní události OnInit, OnStep a OnReset, které řídí logiku prostředí a průběh epizod.
\end{itemize}
\item Trenérský modul:
\begin{itemize}
\item Následuje vytvoření Blueprintu děděného z NevarokMLTrainer, kde se nastavují pozorovací a akční prostory (OnConstruct) a spouští se samotný trénink (OnStart).
\end{itemize}
\item Implementace v úrovni:
\begin{itemize}
\item Trenér i prostředí se umisťují do scény v Unreal Enginu, přičemž trenér odkazuje na prostředí a celý trénink začíná při spuštění úrovně.
\end{itemize}

\newpage

\item Sledování učení:
\begin{itemize}
\item Pro vizualizaci průběhu učení se používá TensorBoard, který umožňuje sledovat vývoj odměn, změny v politice agenta nebo ladění hyperparametrů.
\end{itemize}
\end{enumerate}

Plugin je dostupný pro Unreal Engine ve verzi 5.2 a po jeho aktivaci je nutné spustit i backend klienta a ověřit připojení mezi Unreal Enginem a Pythonovým prostředím.

Hlavní vlastnosti a využití:
\begin{enumerate}
\item Široký výběr algoritmů:
\begin{itemize}
\item Podpora různých typů RL algoritmů umožňuje ladit chování agentů podle konkrétního použití (např. hodnotové nebo policy-gradient přístupy).
\end{itemize}
\item Podpora Blueprintů i C++:
\begin{itemize}
\item Plugin lze plně využívat jak přes vizuální skriptování, tak i přes nativní C++ rozhraní pro pokročilejší zásahy.
\end{itemize}
\item Cross-platform nasazení:
\begin{itemize}
\item I když trénink typicky probíhá na Windows, natrénované modely je možné nasadit na jakoukoli platformu podporovanou Unreal Enginem.
\end{itemize}
\item Vhodné pro výuku a výzkum:
\begin{itemize}
\item Plugin je volně dostupný pro nekomerční využití, což z něj činí ideální nástroj pro akademické projekty, simulace nebo vývoj prototypů.
\end{itemize}
\end{enumerate}

\newpage

\section{Návrh hry}
\subsection{Návrh prostředí hry}
Prostředí, ve kterém bude autonomní agent fungovat, lze realizovat různými způsoby – od jednoduchých po technicky náročnější. Mezi přístupnější varianty patří využití herních enginů, jako je Unity nebo Unreal Engine, zatímco pokročilejší řešení zahrnuje programování vizualizace od základu pomocí nízkoúrovňových API, jako jsou OpenGL nebo DirectX.

V tomto projektu bude vizualizace agenta realizována v Unreal Engine formou střílečky z pohledu první osoby. Prvním krokem bude vytvoření herní mapy, která umožní volný pohyb hráče i postavy řízené umělou inteligencí. Mapa bude inspirována městským náměstím po bitvě – její střed bude tvořit park s bankovní budovou. Park bude téměř prázdný kvůli zpustošení a bankovní budova bude obklopena barikádami. Právě v této budově se hráč po spuštění hry objeví. Bude zde možné nalézt zbraně a box se zbraněmi a brněním, které hráč využije v boji.

Kromě parku bude mapa zahrnovat i zastavěnou oblast – uzavřenější prostor obklopený budovami, který bude představovat vyšší riziko kvůli omezené viditelnosti a pohybu. Uvnitř zabarikádované zóny se budou nacházet výchozí body (spawn pointy) pro nepřátelské postavy ovládané AI.

Pro vyvážení obtížnosti budou nepřátelé po smrti odhazovat sbíratelné předměty, jako jsou lékárničky a zásobníky munice, které hráč využije k doplnění zdraví nebo střeliva.

Součástí uživatelského rozhraní bude tzv. HUD (Head-Up Display), který bude zobrazovat klíčové herní informace – aktuální stav zdraví, zásoby energie, počet nábojů a další. Z tohoto rozhraní bude možné přistupovat jak k hlavnímu menu, tak k pauzovací nabídce.

Celá hra bude strukturována do kolových vln. Po spuštění prvního kola se zavřou dveře banky, přičemž venku bude umístěna další bedna s vybavením, která umožní výměnu zbraní a brnění. V každém kole bude úkolem hráče eliminovat předem daný počet nepřátel. Po jejich porážce se aktuální kolo uzavře a automaticky začne nové – s vyšší obtížností nebo větším počtem protivníků. Tím vznikne dynamický cyklus, který otestuje schopnosti hráče i chování autonomního agenta v různých podmínkách.

Část prostředí a kódu u některých Blueprintů byla převzata z mého předchozího projektu. Vzhledem k rozsáhlým úpravám však bude v následující kapitole podrobně popsána struktura jednotlivých částí hry a následně samotného agenta.

\newpage

\subsection{Návrh agenta}
Cílem je vytvořit agenta řízeného pomocí technik strojového učení, který se dokáže samostatně orientovat v herním prostředí, přiblížit se ke svému cíli (hráči) a efektivně využít střelnou zbraň k jeho eliminaci. Tento agent nebude ovládán předem nadefinovanými skripty, ale své chování se bude učit na základě zkušeností pomocí posilovaného učení. Učící algoritmus bude založen na metodě Proximal Policy Optimization (PPO).

Celkový algoritmus agenta vytvořeného pomocí NevarokML je schematicky znázorněn na obrázku~\ref{obrAlg}. Tento vývojový diagram ilustruje jednotlivé fáze interakce agenta s prostředím — od sběru pozorování, přes předání dat do neuronové sítě, predikci a provedení akce, až po získání odměny, uložení zkušeností a následný trénink modelu.

\begin{figure}[H]
    \centering
    \includegraphics[scale=0.8]{2-2.png}
    \caption{Vývojový diagram algoritmu agenta v rámci NevarokML}
    \label{obrAlg}
\end{figure}

Vzhledem k tomu, že framework NevarokML neumožňuje přímou modifikaci architektury neuronové sítě, jsou všechna zlepšení zaměřena na okolní prvky smyčky — zejména strukturu pozorování, množinu dostupných akcí a definici odměn. Tyto prvky přímo ovlivňují kvalitu politiky, kterou agent získává.

Aby měl agent dostatečný přehled o situaci, budou mu prostřednictvím funkce CollectObservations předávány relevantní informace, na jejichž základě bude činit rozhodnutí. Tyto pozorování mu umožní vyhodnotit aktuální stav prostředí a adekvátně reagovat.

Agent bude mít přístup ke třem základním typům akcí:
\begin{itemize}
\item Pohyb,
\item Rotace,
\item Střelba – za předpokladu splnění podmínek (např. není aktivní cooldown, hráč se nachází v zorném poli apod.).
\end{itemize}

Tyto akce budou implementovány ve funkci ApplyActions, která interpretuje výstup neuronové sítě jak během tréninku, tak i v nasazení.

Jádrem efektivního učení bude systém odměn implementovaný ve funkci HandleRewards, jenž určuje, jak bude agent penalizován nebo odměňován za své chování. Bude například odměňován za správnou rotaci, pohyb směrem k cíli nebo udržování optimální vzdálenosti. Zároveň bude trénink omezen tak, aby agent nebyl motivován pouze maximalizovat počet bodů bez reálného splnění cíle. Veškerá motivace bude přímo vázána na výsledek boje.

Na základě tohoto návrhu se očekává, že agent:
\begin{enumerate}
\item Rozpozná pozici cíle na základě poskytnutých pozorování,
\item Přiblíží se na efektivní vzdálenost,
\item Srovná rotaci tak, aby směřoval k cíli,
\item Zahájí palbu a pokusí se cíl eliminovat.
\end{enumerate}

Tato posloupnost chování by měla emergentně vzniknout díky optimalizaci očekávané kumulativní odměny. Agent si rovněž bude moci po každém kole upravit své vybavení na základě konfigurace, kterou měl hráč v předchozím kole, aby zvýšil svou šanci na úspěch.